\chapter{Effets}
\section{Choisir le type d’effet}
\begin{longtable}{@{}L{3.6cm}L{2.8cm}L{7.2cm}@{}}
\toprule Contrôle & Domaine & Intention\\\midrule\endhead
Texture & $-100$ à $+100$ & Adoucir ou renforcer les petits détails.\\
Clarté & $-100$ à $+100$ & Modifier la présence des structures et le contraste local.\\
Correction du voile & $-100$ à $+100$ & Modifier l’impression de voile et la séparation des tons.\\
Grain & 0 à 100 & Ajouter un effet de grain ; ce n’est pas du débruitage.\\
Flou gaussien & 0 à 20 & Adoucir globalement ; ce n’est pas une simulation complète de profondeur de champ.\\
\bottomrule
\end{longtable}
La référence de ces contrôles est zéro. Une valeur négative de texture ou
de clarté ne correspond pas à une réduction du bruit identique au module dédié.

\section{Contrôler les artefacts}
Examiner les contours, les transitions de ciel et les textures fines après
une forte correction. Un gain de clarté peut rendre le bruit plus visible ;
un excès de lissage peut donner une matière artificielle. Vérifier également
à la taille de livraison, car un détail agréable au fort zoom peut être excessif
en vue ajustée.

\paragraph{Exemple.}
Sur un paysage voilé, commencer par une petite correction du voile, puis
rééquilibrer éventuellement la lumière. Comparer avec \touche{Ctrl + Maj + Y}
et examiner les limites ciel/bâtiments au zoom \touche{Ctrl + 1}.
\gestesreglages
